\documentclass[UTF8]{ctexart}
\usepackage{amsmath}
\usepackage{geometry}
\usepackage{xcolor}
\usepackage{hyperref}
\usepackage{amsfonts}
\usepackage{booktabs}
\geometry{a4paper,scale=0.8}

\title{八股问题}
\author{}
\date{}

\begin{document}

\maketitle
\section{机器学习中的损失函数}

机器学习中的损失函数根据任务类型可以分为以下几类:

\subsection{回归任务损失函数}
\begin{itemize}
    \item \textbf{均方误差(MSE)}: $L = \frac{1}{n}\sum_{i=1}^{n}(y_i - \hat{y}_i)^2$,对异常值敏感
    \item \textbf{平均绝对误差(MAE)}: $L = \frac{1}{n}\sum_{i=1}^{n}|y_i - \hat{y}_i|$,对异常值更鲁棒
    \item \textbf{Huber Loss}: 结合MSE和MAE的优点,在误差较小时使用MSE,误差较大时使用MAE
\end{itemize}

\subsection{分类任务损失函数}
\begin{itemize}
    \item \textbf{交叉熵损失(Cross Entropy)}: 
    \begin{itemize}
        \item 二分类: $L = -\frac{1}{n}\sum_{i=1}^{n}[y_i\log(\hat{y}_i) + (1-y_i)\log(1-\hat{y}_i)]$
        \item 多分类: $L = -\frac{1}{n}\sum_{i=1}^{n}\sum_{c=1}^{C}y_{i,c}\log(\hat{y}_{i,c})$
    \end{itemize}
    \item \textbf{Focal Loss}: $L = -\sum_{i=1}^{n}\alpha_i(1-\hat{y}_i)^{\gamma}y_i\log(\hat{y}_i)$,解决类别不平衡
    \item \textbf{KL散度}: $D_{KL}(P||Q) = \sum_{i}P(i)\log\frac{P(i)}{Q(i)}$,衡量两个分布的差异
\end{itemize}

\subsection{对比学习损失函数}
\begin{itemize}
    \item \textbf{InfoNCE Loss}: $L = -\log\frac{\exp(sim(z_i, z_i^+)/\tau)}{\sum_{j=1}^{N}\exp(sim(z_i, z_j)/\tau)}$
\end{itemize}

\section{交叉熵作为损失函数的原理及相关概念}

\subsection{为什么交叉熵能作为损失函数}

交叉熵能作为损失函数的核心原因有以下几点:

\textbf{1. 最大似然估计的等价性}

最小化交叉熵等价于\textbf{最大化似然函数}。假设样本独立同分布,似然函数为:
$$L(\theta) = \prod_{i=1}^{n}P(y_i|x_i;\theta)$$

取负对数得到:
$$-\log L(\theta) = -\sum_{i=1}^{n}\log P(y_i|x_i;\theta)$$

这正是交叉熵损失的形式。所以\textbf{最小化交叉熵就是最大化似然},有坚实的统计学基础。

\textbf{2. 梯度性质良好}

交叉熵配合softmax激活函数,梯度形式简洁:
$$\frac{\partial L}{\partial z_i} = \hat{y}_i - y_i$$

梯度大小正好等于预测误差,不会出现梯度消失问题,训练稳定高效。

\textbf{3. 凸优化性质}

对于线性模型,交叉熵是凸函数,保证能找到全局最优解。

\subsection{InfoNCE与交叉熵的关系}

InfoNCE(Noise Contrastive Estimation)本质上是\textbf{多分类交叉熵的一种特殊形式}。

InfoNCE的形式:
$$L_{InfoNCE} = -\log\frac{\exp(sim(z_i, z_i^+)/\tau)}{\sum_{j=1}^{N}\exp(sim(z_i, z_j)/\tau)}$$

可以看出,这就是一个\textbf{N分类的交叉熵}:
\begin{itemize}
    \item 正样本$z_i^+$对应真实类别,标签为1
    \item 其他N-1个负样本对应其他类别,标签为0
    \item 分母是softmax归一化
\end{itemize}

所以InfoNCE = \textbf{把对比学习转化为N分类问题},用交叉熵优化。核心思想是\textbf{最大化正样本对的相似度,同时最小化负样本对的相似度}。

\subsection{KL散度与交叉熵的关系}

KL散度和交叉熵有紧密的数学关系:

$$D_{KL}(P||Q) = \sum_{i}P(i)\log\frac{P(i)}{Q(i)} = \sum_{i}P(i)\log P(i) - \sum_{i}P(i)\log Q(i)$$

$$D_{KL}(P||Q) = -H(P) + H(P,Q)$$

其中$H(P)$是P的熵,$H(P,Q)$是交叉熵。

\textbf{关键结论}:
\begin{itemize}
    \item KL散度 = 交叉熵 - 熵
    \item 在分类任务中,真实标签P是固定的one-hot分布,熵$H(P)=0$
    \item 因此\textbf{最小化交叉熵等价于最小化KL散度}
    \item KL散度非负,当且仅当P=Q时为0
\end{itemize}

所以在分类任务中,\textbf{最小化交叉熵就是让模型预测分布Q尽可能接近真实分布P},这正是我们想要的。

\textbf{区别}:
\begin{itemize}
    \item KL散度不对称: $D_{KL}(P||Q) \neq D_{KL}(Q||P)$
    \item 交叉熵也不对称,但在分类任务中P固定,只优化Q
    \item KL散度更多用于分布匹配(如VAE),交叉熵更多用于分类任务
\end{itemize}


\section{PPO算法详解及其在大模型后训练中的应用}

\subsection{PPO算法原理}

PPO(Proximal Policy Optimization)是\textbf{目前最流行的强化学习算法},由OpenAI在2017年提出。

\textbf{核心思想}: 在策略更新时,限制新策略和旧策略的差异不要太大,避免训练不稳定。

\textbf{目标函数}:
$$L^{CLIP}(\theta) = \mathbb{E}_t[\min(r_t(\theta)\hat{A}_t, \text{clip}(r_t(\theta), 1-\epsilon, 1+\epsilon)\hat{A}_t)]$$

其中:
\begin{itemize}
    \item $r_t(\theta) = \frac{\pi_\theta(a_t|s_t)}{\pi_{\theta_{old}}(a_t|s_t)}$是重要性采样比率
    \item $\hat{A}_t$是优势函数估计
    \item $\epsilon$是裁剪参数,通常取0.1或0.2
\end{itemize}

\textbf{关键机制}:
\begin{enumerate}
    \item \textbf{裁剪(Clipping)}: 当$r_t$超出$[1-\epsilon, 1+\epsilon]$范围时,梯度被截断,防止策略更新过大
    \item \textbf{优势函数}: 用GAE(Generalized Advantage Estimation)估计,减小方差
    \item \textbf{多轮更新}: 同一批数据可以重复使用多次,提高样本效率
\end{enumerate}

\textbf{算法流程}:
\begin{enumerate}
    \item 用当前策略$\pi_{\theta_{old}}$收集轨迹数据
    \item 计算每个时间步的优势函数$\hat{A}_t$
    \item 用收集的数据更新策略,最大化$L^{CLIP}$
    \item 重复K轮(通常K=3-10)
    \item 更新$\theta_{old} \leftarrow \theta$,进入下一轮采样
\end{enumerate}

\subsection{为什么大模型后训练选择PPO}

大模型的RLHF(Reinforcement Learning from Human Feedback)后训练几乎都选择PPO,原因如下:
\begin{itemize}
    \item[1] 训练稳定性：大模型参数量巨大(几十亿到千亿),训练极其不稳定，PPO的\textbf{裁剪机制}保证策略更新幅度可控,避免灾难性遗忘。
    \item[2] 工程实现友好：相比TRPO,PPO实现简单,不需要计算二阶导数和共轭梯度；支持\textbf{批量并行采样},充分利用GPU集群
    \item[3] 样本效率：大模型生成样本成本极高(需要推理生成文本)，PPO是\textbf{on-policy算法},但允许\textbf{多轮更新}同一批数据，实现了稳定性与效率的平衡
    \item[4] 超参数鲁棒性：大模型训练成本高,不能频繁调参，而PPO对超参数不敏感
\end{itemize}


\subsection{RLHF中的PPO应用}

在ChatGPT等大模型的RLHF训练中,PPO的具体应用流程:
\begin{itemize}
    \item[1] 阶段1: 监督微调(SFT)————用高质量人工标注数据微调预训练模型，得到初始策略$\pi_{SFT}$
    \item[2] 阶段2: 奖励模型训练(RM)————人工对模型生成的多个回复进行排序，训练奖励模型$r_\phi(x,y)$预测人类偏好
    \item[3] 阶段3: PPO强化学习————用户输入prompt $x$，模型生成回复$y \sim \pi_\theta(y|x)$，计算奖励$r(x,y) = r_\phi(x,y) - \beta \log\frac{\pi_\theta(y|x)}{\pi_{SFT}(y|x)}$（第二项是KL惩罚,防止模型偏离SFT太远），用PPO最大化期望奖励
\end{itemize}

总结: PPO在\textbf{稳定性、样本效率、实现简单性}之间达到了最佳平衡,是大模型RLHF的不二之选。

\section{MCP和Skill的定义与联系}
MCP和Skill是两种不同的能力扩展机制，一句话总结：「MCP 解决”连接”问题：让 AI 能访问外部世界」「Skill 解决”方法论”问题：教 AI 怎么做某类任务」
\begin{itemize}
\item[1]MCP定义：MCP是一种标准化协议,允许AI系统通过统一接口调用外部工具和服务（官方的比喻是”AI 应用的 type-C 接口”——就像 type-C 提供了一种通用的方式连接各种设备，MCP 提供了一种通用的方式连接各种工具和数据源。）
\item[2]为什么需要MCP：大模型本身\textbf{无法直接操作外部系统}(数据库、文件系统、API),MCP提供了\textbf{行动能力}；其次大模型的知识截止于训练时间,MCP可以\textbf{获取最新数据}；大模型在复杂计算上容易出错,MCP可以调用\textbf{专业工具}(计算器、代码执行器)；无需修改AI核心,通过\textbf{添加MCP服务器}即可扩展新能力
\item[3]Skill定义：Skill是预定义的\textbf{指令集或提示模板},增强AI在特定领域的能力，直接\textbf{注入到AI的上下文}中,成为系统提示的一部分，所谓“一次实现，到处使用”
\item[4]为什么需要Skill：首先，Skill可以注入领域专家的经验和最佳实践使得大模型在指定领域表现卓越；其次，Skill可以\textbf{引导AI的推理方向},通过明确的指令和约束，避免偏离任务目标或产生不相关的输出；并且，大模型的训练数据有时效性,Skill可以\textbf{补充最新知识}
\item[5]两者通常\textbf{配合使用}
\end{itemize}

\textbf{Skill + MCP的协同价值}

\begin{itemize}
    \item \textbf{知识与行动结合}: Skill提供\textbf{"知道怎么做"},MCP提供\textbf{"能够做"}
    \item \textbf{完整的工作流}: Skill定义流程,MCP执行具体操作,形成闭环
    \item \textbf{质量保证}: Skill确保\textbf{决策质量},MCP确保\textbf{执行准确性}
\end{itemize}

\textbf{总结}: Skill和MCP分别解决了AI的\textbf{"智力"和"能力"}问题。Skill让AI\textbf{智力更高}(知道该做什么、怎么做好),MCP让AI\textbf{能力更强}(能够操作真实世界)。两者结合,才能构建真正\textbf{实用的AI系统}。

\section{传统强化学习Agent与大模型Agent的区别}

\subsection{传统强化学习Agent}

\textbf{核心特征}:
\begin{itemize}
    \item \textbf{学习方式}: 通过\textbf{试错}与环境交互,从奖励信号中学习策略
    \item \textbf{状态表示}: 通常是\textbf{低维向量}或离散状态(如游戏画面、机器人传感器数据)
    \item \textbf{动作空间}: 明确定义的\textbf{有限动作集}(如上下左右、关节角度)
    \item \textbf{策略形式}: 神经网络映射状态到动作概率分布
    \item \textbf{优化目标}: 最大化\textbf{累积奖励}
    \item \textbf{泛化能力}: 局限于\textbf{训练环境},难以迁移到新任务
\end{itemize}

\textbf{典型应用}: AlphaGo(围棋)、Atari游戏、机器人控制

\subsection{大模型Agent}

\textbf{核心特征}:
\begin{itemize}
    \item \textbf{学习方式}: 基于\textbf{预训练+微调},从海量文本中学习世界知识和推理能力
    \item \textbf{状态表示}: \textbf{自然语言描述}的复杂场景,可以理解抽象概念
    \item \textbf{动作空间}: \textbf{开放式}动作,可以生成任意文本、调用工具、执行代码
    \item \textbf{策略形式}: 语言模型生成下一步行动的自然语言描述
    \item \textbf{优化目标}: 完成\textbf{用户指定的任务目标},可以是多步推理、规划、执行
    \item \textbf{泛化能力}: 强大的\textbf{零样本/少样本}学习能力,可以处理未见过的任务
\end{itemize}

\textbf{典型应用}: ChatGPT、AutoGPT、代码助手、智能客服

\subsection{二者不同}
\begin{itemize}
\item[(1)]训练方式不同：传统agent能力来自于交互中的试错；而大模型agent依赖于预训练得到的语言知识 + 提示词 + 工具调用 + 外部记忆
\item[(2)]环境类型不同：传统agent适用于状态清晰、动作可枚举、奖励可定义的问题；而大模型agent适用于开放世界、非结构化、语言主导的任务
\item[(3)]动作空间不同：传统agent动作通常是离散或连续控制信号，而大模型agent的动作更像高层语义动作
\item[(4)]总结：传统强化学习智能体是“以奖励驱动的策略学习器”，大模型智能体是“以大模型为中枢的任务执行系统”；前者重点在学策略，后者重点在组织知识、规划步骤、调用工具和完成复杂开放任务。
\end{itemize}
\section{AI Coding的Plan模式与Agent模式区别}

在AI Coding系统中，Plan模式与Agent模式代表了两种不同的任务组织方式。前者强调先形成较完整的实现方案，再按既定思路执行；后者则强调在执行过程中持续获取反馈，并根据中间结果动态调整后续动作。

\subsection{Plan模式}

Plan模式的核心是“先规划，后执行”。当用户提出任务后，系统先分析需求，明确需要修改的文件、涉及的函数以及大致实施步骤，在此基础上再生成或修改代码。其优势在于流程清晰、成本较低、执行效率较高，因而更适合目标明确、影响范围有限的任务，例如局部代码生成、简单重构、变量重命名以及格式化修复等。

不过，Plan模式对前期规划质量较为依赖。如果任务本身存在较强不确定性，或者执行过程中需要频繁依赖编译结果、测试反馈和运行日志来修正判断，那么预先给出的方案可能很快失效，因此其应对复杂场景的能力相对有限。

\subsection{Agent模式}

Agent模式的核心是“边执行，边调整”。系统通常先观察当前代码库状态、错误信息或测试结果，再决定下一步行动，例如读取文件、搜索符号、修改代码、运行测试或调用外部工具；随后根据新的反馈继续判断和修正，直至任务完成。其本质是一种“观察—行动—反馈—再决策”的闭环过程。

相较于Plan模式，Agent模式更适合复杂缺陷修复、多文件协同修改、需要调试验证的功能实现以及与外部系统交互的任务。其优势在于能够处理不确定性并具备较强自适应能力，但代价是推理轮次更多、工具调用更频繁，因此整体成本和时延通常也更高。

\subsection{两种模式的区别与实际应用}

总体而言，Plan模式更强调执行前的集中规划，适合简单、明确、可控的任务；Agent模式更强调执行中的动态决策，适合复杂、开放、依赖反馈的任务。前者通常更快、更省成本，但对异常情况较为敏感；后者通常更稳健、更灵活，但系统复杂度和资源开销也更高。

在实际AI Coding系统中，二者往往并非完全割裂，而是以混合方式共同存在。对于多数中高复杂度的AI Coding任务，更有效的使用方式通常不是单独采用Plan模式或Agent模式，而是先由plan生成可审阅的任务计划，以明确整体修改方向、影响范围与验证路径，再由Agent在执行过程中结合编译结果、测试反馈与环境状态进行动态调整。对于简单且边界清晰的小规模任务，则可直接进入Agent执行阶段，以避免过度规划带来的额外开销。

\section{Vibe Coding与Spec Coding的区别}

Vibe Coding和Spec Coding代表了AI辅助编程中两种不同的交互范式，前者强调即时性和流畅性，后者强调结构化和可控性。

\subsection{Vibe Coding}

Vibe Coding的核心是\textbf{"跟着感觉走，快速迭代"}。开发者通过自然语言与AI进行\textbf{对话式交互}，描述想法或需求，AI立即生成代码或进行修改，开发者查看结果后继续提出调整意见，形成快速的反馈循环。整个过程类似于与一位结对编程伙伴的即兴协作，强调\textbf{思维流畅性}和\textbf{创造性探索}。

\textbf{典型特征}:
\begin{itemize}
    \item \textbf{交互方式}: 自由对话，需求可以模糊、渐进式明确
    \item \textbf{工作流程}: 提出想法 → AI生成 → 查看结果 → 继续调整（循环）
    \item \textbf{适用场景}: 原型开发、探索性编程、快速验证想法、学习新技术
    \item \textbf{心智模型}: "边想边做"，允许需求在过程中演化
\end{itemize}

\textbf{优势}:
\begin{itemize}
    \item 上手门槛低，无需预先完整规划
    \item 快速启动，适合灵感驱动的开发
    \item 灵活性高，可随时调整方向
    \item 学习曲线平缓，适合初学者
\end{itemize}

\textbf{劣势}:
\begin{itemize}
    \item 缺乏整体规划，容易产生技术债务
    \item 对话历史过长时，AI可能遗忘早期上下文
    \item 难以追溯决策过程和变更历史
    \item 复杂项目中容易失控，代码质量参差不齐
\end{itemize}

\subsection{Spec Coding}

Spec Coding的核心是\textbf{"先定义规格，再分步实现"}。开发者首先编写结构化的\textbf{规格文档}(Spec)，明确需求、设计方案、实现任务和验收标准，然后AI根据Spec逐步执行任务。整个过程类似于传统软件工程中的需求分析和设计阶段，强调\textbf{可控性}和\textbf{可追溯性}。

\textbf{典型特征}:
\begin{itemize}
    \item \textbf{交互方式}: 结构化文档，需求需要预先明确
    \item \textbf{工作流程}: 编写Spec → 审阅计划 → AI执行任务 → 验证结果
    \item \textbf{适用场景}: 复杂功能开发、多人协作、需要文档化的项目、关键系统
    \item \textbf{心智模型}: "先想清楚再做"，需求相对稳定
\end{itemize}

\textbf{优势}:
\begin{itemize}
    \item 整体规划清晰，代码质量更可控
    \item 任务可拆解、可并行、可追溯
    \item 适合团队协作，Spec可作为沟通文档
    \item 支持渐进式开发，可分阶段验证
    \item 便于代码审查和知识传承
\end{itemize}

\textbf{劣势}:
\begin{itemize}
    \item 前期投入较大，需要编写详细Spec
    \item 灵活性较低，需求变更需要更新Spec
    \item 学习成本较高，需要掌握Spec编写规范
    \item 对于简单任务可能过度设计
\end{itemize}

\subsection{两种模式的对比与选择}

\begin{table}[h]
\centering
\begin{tabular}{lll}
\toprule
\textbf{维度} & \textbf{Vibe Coding} & \textbf{Spec Coding} \\
\midrule
交互方式 & 自由对话 & 结构化文档 \\
需求明确度 & 可模糊、渐进式 & 需预先明确 \\
启动速度 & 快 & 慢 \\
适用复杂度 & 简单到中等 & 中等到复杂 \\
代码质量 & 不稳定 & 较稳定 \\
可追溯性 & 弱 & 强 \\
团队协作 & 困难 & 友好 \\
学习曲线 & 平缓 & 陡峭 \\
\bottomrule
\end{tabular}
\end{table}

\textbf{实际应用建议}:
\begin{itemize}
    \item \textbf{快速原型、探索性开发}: 优先使用Vibe Coding
    \item \textbf{生产级功能、复杂系统}: 优先使用Spec Coding
    \item \textbf{混合使用}: 用Vibe Coding快速验证想法，确定方向后转为Spec Coding进行正式开发
    \item \textbf{个人项目}: 根据个人偏好选择，Vibe Coding更自由
    \item \textbf{团队项目}: Spec Coding更适合协作和知识共享
\end{itemize}

\textbf{总结}: Vibe Coding强调\textbf{速度和灵活性}，适合探索和快速迭代；Spec Coding强调\textbf{质量和可控性}，适合复杂项目和团队协作。两者并非对立，而是互补关系，开发者应根据项目阶段、复杂度和团队情况灵活选择。

\end{document}
